\chapter{Bayesian networks and d-separation}
\label{app:bayesnet}

The derivations of \cref{ch:statespace,ch:bayes} repeatedly drop variables
from a conditioning: ``given the current pose, the past does not matter'',
``the next control says nothing about the current pose''. Each of these is
a statement of \emph{conditional independence}. This appendix explains a
way to check such statements without any computation: draw the model as a
graph, and read the independences off the graph.

It proceeds in four steps: what conditional independence means
(\cref{sec:condind}); how a model is drawn as a graph
(\cref{sec:bn}); the three situations that can occur along a path of the
graph, each proved from the definition (\cref{sec:three}); and the
general rule, called \emph{d-separation}, with a recipe to apply it
(\cref{sec:dsep}). \Cref{sec:bayesnet-mower} then applies the recipe to
the mower.

\section{Conditional independence}
\label{sec:condind}

Let $A$, $B$ and $C$ be random variables. $A$ and $B$ are
\emph{independent} if knowing one says nothing about the other:
$p(a\mid b)=p(a)$, or equivalently $p(a,b)=p(a)\,p(b)$. We write
$A\indep B$.

They are \emph{conditionally independent given $C$} if, \emph{once $C$ is
known}, knowing $B$ says nothing more about $A$:
\[
  p(a\mid b,c)=p(a\mid c),\qquad\text{or equivalently}\qquad
  p(a,b\mid c)=p(a\mid c)\,p(b\mid c).
\]
We write $A\indep B\mid C$. This is exactly what allows a variable to be
dropped from a conditioning.

The two notions are different, and neither implies the other. Two
variables can be dependent and become independent once a third is known;
two variables can be independent and become dependent once a third is
known. \Cref{sec:three} gives an example of each.

\section{Bayesian networks}
\label{sec:bn}

A model is often described by saying, for each variable, which other
variables it is computed from. A \emph{Bayesian network} draws this
description: one node per variable, and an arrow $A\rightarrow B$ when $A$
is one of the variables that $B$ is computed from. The arrows must not
form a cycle.

Some vocabulary. In $A\rightarrow B$, $A$ is a \emph{parent} of $B$ and
$B$ a \emph{child} of $A$. The set of parents of a node $v$ is written
$\mathrm{pa}(v)$. The \emph{descendants} of $v$ are its children, their
children, and so on: every node that can be reached from $v$ by following
arrows forwards.

The graph stands for a factorisation of the joint distribution, with one
factor per variable, the distribution of that variable given its parents:
\begin{equation}\label{eq:bnfactor}
  p(v_1,\dots,v_n)=\prod_{k=1}^n p\bigl(v_k\mid\mathrm{pa}(v_k)\bigr).
\end{equation}
A variable without parents contributes its plain distribution $p(v_k)$.

\begin{example}{A garden}{garden}
Four variables: it rained ($R$), the sprinkler was on ($S$), the grass is
wet ($W$), someone slips on the grass ($L$). Rain and sprinkler each make
the grass wet, and wet grass makes people slip (\cref{fig:garden}). The
parents are $\mathrm{pa}(W)=\{R,S\}$ and $\mathrm{pa}(L)=\{W\}$; $R$ and
$S$ have none. The descendants of $R$ are $W$ and $L$. The factorisation
\eqref{eq:bnfactor} reads
\[
  p(r,s,w,l)=p(r)\,p(s)\,p(w\mid r,s)\,p(l\mid w).
\]
\end{example}

\begin{figure}[ht]
\centering
\begin{tikzpicture}[>=Stealth,
  var/.style={circle,draw,thick,minimum size=8mm,inner sep=0pt,font=\small}]
  \node[var] (R) at (0,1.4) {$R$};
  \node[var] (S) at (3,1.4) {$S$};
  \node[var] (W) at (1.5,0) {$W$};
  \node[var] (L) at (1.5,-1.5) {$L$};
  \draw[->,thick] (R) -- (W);
  \draw[->,thick] (S) -- (W);
  \draw[->,thick] (W) -- (L);
  \node[left=2pt,font=\small] at (R.west) {rain};
  \node[right=2pt,font=\small] at (S.east) {sprinkler};
  \node[right=2pt,font=\small] at (W.east) {wet grass};
  \node[right=2pt,font=\small] at (L.east) {slip};
\end{tikzpicture}
\caption{The Bayesian network of \cref{ex:garden}.}
\label{fig:garden}
\end{figure}

For the mower (\cref{fig:bayesnet}), the pose $\x_{t+1}$ is computed from
$\x_t$ and $\uu_t$, and the measurement $\z_t$ from $\x_t$:
$\mathrm{pa}(\x_{t+1})=\{\x_t,\uu_t\}$ and $\mathrm{pa}(\z_t)=\{\x_t\}$.
The factorisation \eqref{eq:bnfactor} is then \eqref{eq:joint}. The
controls are known inputs rather than random variables; they are drawn as
nodes that are always given, and nothing below changes.

\section{The three building blocks}
\label{sec:three}

Take three variables $A$, $B$, $C$ with an arrow between $A$ and $B$ and
an arrow between $B$ and $C$, and no arrow between $A$ and $C$. The two
arrows can be oriented in three essentially different ways
(\cref{fig:connections}). In each case the question is the same: are $A$
and $C$ independent, and does the answer change when $B$ is known?

\begin{figure}[ht]
\centering
\begin{tikzpicture}[>=Stealth,
  var/.style={circle,draw,thick,minimum size=7mm,inner sep=0pt,font=\small,
    fill=white},
  giv/.style={fill=defblue!22}]
  \node[font=\small\bfseries] at (2.6,1.0) {$B$ not given};
  \node[font=\small\bfseries] at (8.6,1.0) {$B$ given};
  \foreach \row/\name/\ta/\tb/\va/\vb in {
      0/chain/->/->/open/blocked,
      -1.5/fork/<-/->/open/blocked,
      -3.0/collider/->/<-/blocked/open}{
    \node[font=\small,anchor=east] at (0.6,\row) {\name};
    \foreach \dx/\st/\verdict in {1.2/{}/\va, 7.2/giv/\vb}{
      \node[var] (A) at (\dx,\row) {$A$};
      \node[var,\st] (B) at (\dx+1.4,\row) {$B$};
      \node[var] (C) at (\dx+2.8,\row) {$C$};
      \draw[\ta,thick] (A) -- (B);
      \draw[\tb,thick] (B) -- (C);
      \node[font=\small,anchor=west] at (\dx+3.4,\row) {\verdict};
    }
  }
\end{tikzpicture}
\caption{The three ways two arrows can meet at a middle node $B$, and
whether information passes between $A$ and $C$ (``open'') or not
(``blocked''). A shaded node is given. Chains and forks are open until $B$
is given; a collider is the other way round.}
\label{fig:connections}
\end{figure}

\subsection*{Chain: $A\rightarrow B\rightarrow C$}

\emph{Example.} It rains ($A$), so the grass is wet ($B$), so someone slips
($C$). Rain makes slipping more likely: $A$ and $C$ are dependent. But if
we already know whether the grass is wet, learning that it rained tells
nothing more about slipping, because rain acts on slipping only through
the wet grass.

\emph{Statement.} In a chain, $A$ and $C$ are in general dependent, and
$A\indep C\mid B$.

\emph{Proof.} The factorisation \eqref{eq:bnfactor} is
$p(a,b,c)=p(a)\,p(b\mid a)\,p(c\mid b)$. Its first two factors multiply to
$p(a,b)$, so
\[
  p(c\mid a,b)=\frac{p(a,b,c)}{p(a,b)}=p(c\mid b).
\]
Given $B$, the value of $A$ does not change the distribution of $C$.

\subsection*{Fork: $A\leftarrow B\rightarrow C$}

\emph{Example.} A patient has the flu ($B$), which causes fever ($A$) and
cough ($C$). In the population, fever and cough go together: seeing one
makes the other more likely, because it makes the flu more likely. Among
patients known to have the flu, a fever says nothing more about a cough.

\emph{Statement.} In a fork, $A$ and $C$ are in general dependent, and
$A\indep C\mid B$.

\emph{Proof.} The factorisation is $p(a,b,c)=p(b)\,p(a\mid b)\,p(c\mid b)$.
Dividing by $p(b)$,
\[
  p(a,c\mid b)=p(a\mid b)\,p(c\mid b).
\]

In both cases the given middle node is said to \emph{screen off} $A$ from
$C$.

\subsection*{Collider: $A\rightarrow B\leftarrow C$}

\emph{Example.} Rain ($A$) and the sprinkler ($C$) both wet the grass
($B$). Rain and sprinkler have nothing to do with each other: they are
independent. Now suppose we see that the grass is wet. If we then learn
that the sprinkler was on, rain becomes \emph{less} likely, because the
sprinkler already explains the wet grass. Once the common effect is known,
its two causes are no longer independent. This is called \emph{explaining
away}.

\emph{Statement.} In a collider, $A\indep C$, and in general $A$ and $C$
are dependent given $B$.

\emph{Proof.} The factorisation is $p(a,b,c)=p(a)\,p(c)\,p(b\mid a,c)$.
Summing over $b$, and using $\sum_bp(b\mid a,c)=1$,
\[
  p(a,c)=p(a)\,p(c):
\]
$A$ and $C$ are independent. Given $B$, however,
$p(a,c\mid b)=p(a)\,p(c)\,p(b\mid a,c)/p(b)$, and the factor
$p(b\mid a,c)$ ties $a$ and $c$ together; the next example shows that the
dependence is real.

\begin{example}{Explaining away, with numbers}{explain}
It rains with probability $0.2$, the sprinkler is on with probability
$0.2$, independently, and the grass is wet exactly when at least one of
the two happened. The grass is wet with probability $1-0.8\times0.8=0.36$.
\begin{itemize}[nosep]
  \item Knowing only that the grass is wet, the probability of rain is
    $0.2/0.36=0.56$.
  \item Knowing that the grass is wet \emph{and} that the sprinkler was on,
    the grass would be wet whether it rained or not, so the wet grass says
    nothing about rain: its probability is back to $0.2$.
\end{itemize}
Given the wet grass, learning about the sprinkler changes the probability
of rain from $0.56$ to $0.2$: the two causes are dependent.
\end{example}

One more fact about colliders: knowing a \emph{descendant} of the collider
has the same effect as knowing the collider. If we do not see the grass
but see someone slip on it, the grass is probably wet, and rain and
sprinkler become dependent again.

\begin{inwords}
Information flows through a chain or a fork, unless the middle node is
known. Information does not flow through a collider, unless the collider,
or something it causes, is known.
\end{inwords}

\section{d-separation}
\label{sec:dsep}

In a larger graph, two variables can be connected by several routes, each
made of chains, forks and colliders. The rule of d-separation (``d'' for
\emph{directional}) combines the three building blocks.

A \emph{path} between two nodes is a sequence of nodes in which each one is
joined to the next by an arrow, \emph{whatever the direction of the
arrows}. In \cref{fig:garden}, $R\rightarrow W\leftarrow S$ is a path from
$R$ to $S$, although no arrow leads from $R$ to $S$. Each intermediate node
of a path is a chain, a fork or a collider \emph{on that path}, according
to the two arrows of the path that meet at it.

\begin{definition}{Blocked path, d-separation}{dsep}
Let $\mathcal S$ be a set of given variables. A path is \emph{blocked} by
$\mathcal S$ if at least one of its intermediate nodes $B$ satisfies one
of the following:
\begin{itemize}[nosep]
  \item $B$ is a chain or a fork on the path, and $B$ is in $\mathcal S$;
  \item $B$ is a collider on the path, and neither $B$ nor any descendant
    of $B$ is in $\mathcal S$.
\end{itemize}
A path that is not blocked is \emph{open}. Two variables $X$ and $Y$ are
\emph{d-separated} by $\mathcal S$ if every path between them is blocked
by $\mathcal S$.
\end{definition}

One blocking node is enough to block a path, but every path must be
blocked for d-separation.

\begin{theorem}{d-separation implies conditional independence}{dsep}
If $X$ and $Y$ are d-separated by $\mathcal S$ in a Bayesian network, then
$X\indep Y\mid\mathcal S$ for every distribution that factorises as
\eqref{eq:bnfactor}.
\end{theorem}

The proofs of \cref{sec:three} are the three smallest cases of this
theorem. The general proof extends the same computations to any graph;
it can be found in Koller and Friedman (see the references).

\paragraph{A recipe.} To decide whether $X\indep Y\mid\mathcal S$:
\begin{enumerate}[nosep]
  \item Shade the nodes of $\mathcal S$.
  \item List the paths between $X$ and $Y$, ignoring the direction of the
    arrows.
  \item On each path, look for a blocking node: a shaded chain or fork
    node, or a collider that is unshaded and has no shaded descendant.
  \item If every path has a blocking node, $X$ and $Y$ are independent
    given $\mathcal S$. If some path has none, the graph does not
    guarantee independence.
\end{enumerate}

\begin{example}{Three questions about the garden}{gardenq}
In \cref{fig:garden}:
\begin{itemize}[nosep]
  \item \emph{Is $R\indep S$, with nothing given?} The only path is
    $R\rightarrow W\leftarrow S$. Its middle node $W$ is a collider; it is
    not given, and neither is its descendant $L$. The path is blocked:
    yes, rain and sprinkler are independent.
  \item \emph{Is $R\indep S\mid L$?} Same path, but now $L$, a descendant
    of the collider $W$, is given. The path is open: no.
  \item \emph{Is $R\indep L\mid W$?} The only path is
    $R\rightarrow W\rightarrow L$, a chain whose middle node $W$ is given.
    The path is blocked: yes. Once the state of the grass is known, rain
    says nothing more about slipping.
\end{itemize}
\end{example}

\section{Reading the mower's network}
\label{sec:bayesnet-mower}

The recipe is now applied to the four statements that the book relies on.
Each is drawn on four consecutive steps of the network of
\cref{fig:bayesnet}; the rest of the chain, to the left and to the right,
only adds paths that pass through the same nodes. In each drawing the
shaded nodes are given, the nodes with an orange ring are the two whose
independence is in question, and a coloured band marks a path between
them: red if it is blocked, green if it is open.

Two features of this network do most of the work. First, a measurement
$\z_t$ has a single neighbour, its parent $\x_t$. Second, the poses form a
chain, so that the past and the future are joined only through the
current pose.

\subsection*{(a) Assumption \ref{A2}: the measurement depends only on the
current pose}

\emph{Claim.} $p(\z_t\mid\x_{0:t},\uu_{0:t-1},\z_{1:t-1})=p(\z_t\mid\x_t)$:
given $\x_t$, the measurement $\z_t$ is independent of all the earlier
poses, controls and measurements.

\begin{center}
\begin{tikzpicture}[scale=1.25,>=Stealth,
  var/.style={circle,draw,thick,minimum size=9.5mm,inner sep=0pt,fill=white,font=\small},
  ctl/.style={rectangle,draw,minimum size=6.5mm,inner sep=1.5pt,fill=white,font=\small},
  giv/.style={fill=defblue!22},
  qry/.style={draw=orange!85!black,line width=1.6pt}]
  \draw[line width=5pt,red!70,opacity=0.35,line cap=round,line join=round] (0.000,-1.250) -- (0.000,0.000) -- (1.650,0.000) -- (1.650,-1.250);
  \draw[line width=5pt,red!70,opacity=0.35,line cap=round,line join=round] (0.825,1.050) -- (1.650,0.000) -- (1.650,-1.250);
  \node[var,giv] (x0) at (0.000,0.000) {$\x_{t-1}$};
  \node[var,giv] (x1) at (1.650,0.000) {$\x_{t}$};
  \node[var] (x2) at (3.300,0.000) {$\x_{t+1}$};
  \node[var] (x3) at (4.950,0.000) {$\x_{t+2}$};
  \node[var,giv] (z0) at (0.000,-1.250) {$\z_{t-1}$};
  \node[var,qry] (z1) at (1.650,-1.250) {$\z_{t}$};
  \node[var] (z2) at (3.300,-1.250) {$\z_{t+1}$};
  \node[var] (z3) at (4.950,-1.250) {$\z_{t+2}$};
  \node[ctl,giv] (u0) at (0.825,1.050) {$\uu_{t-1}$};
  \node[ctl] (u1) at (2.475,1.050) {$\uu_{t}$};
  \node[ctl] (u2) at (4.125,1.050) {$\uu_{t+1}$};
  \draw[->,thick] (x0) -- (x1);
  \draw[->,thick] (x1) -- (x2);
  \draw[->,thick] (x2) -- (x3);
  \draw[->,thick] (u0) -- (x1);
  \draw[->,thick] (u1) -- (x2);
  \draw[->,thick] (u2) -- (x3);
  \draw[->,thick] (x0) -- (z0);
  \draw[->,thick] (x1) -- (z1);
  \draw[->,thick] (x2) -- (z2);
  \draw[->,thick] (x3) -- (z3);
  \draw[->,thick,dotted] ($(x0)+(-1.1,0)$) -- (x0);
  \draw[->,thick,dotted] (x3) -- ++(1.1,0);
\end{tikzpicture}

\end{center}

\emph{Check.} Every path that reaches $\z_t$ ends with the arrow
$\x_t\rightarrow\z_t$, because $\x_t$ is its only neighbour. On such a
path $\x_t$ is an intermediate node with an arrow leaving towards $\z_t$,
so it is a chain node (if the path arrives from $\x_{t-1}$ or
$\uu_{t-1}$) or a fork node (if it arrives from $\x_{t+1}$), never a
collider. It is given, so every path is blocked. The drawing shows two of
them, from $\z_{t-1}$ and from $\uu_{t-1}$.

\subsection*{(b) Assumption \ref{A1}: the next pose depends only on the
current pose and control}

\emph{Claim.} $p(\x_{t+1}\mid\x_{0:t},\uu_{0:t},\z_{1:t})=
p(\x_{t+1}\mid\x_t,\uu_t)$: given $\x_t$ and $\uu_t$, the next pose is
independent of the earlier poses $\x_{0:t-1}$, the earlier controls
$\uu_{0:t-1}$ and the measurements $\z_{1:t}$.

\begin{center}
\begin{tikzpicture}[scale=1.25,>=Stealth,
  var/.style={circle,draw,thick,minimum size=9.5mm,inner sep=0pt,fill=white,font=\small},
  ctl/.style={rectangle,draw,minimum size=6.5mm,inner sep=1.5pt,fill=white,font=\small},
  giv/.style={fill=defblue!22},
  qry/.style={draw=orange!85!black,line width=1.6pt},
  drp/.style={draw=violet!80!black,line width=1.6pt}]
  \draw[line width=5pt,red!70,opacity=0.35,line cap=round,line join=round] (0.000,-1.250) -- (0.000,0.000) -- (1.650,0.000) -- (3.300,0.000);
  \draw[line width=5pt,red!70,opacity=0.35,line cap=round,line join=round] (1.650,-1.250) -- (1.650,0.000) -- (3.300,0.000);
  \draw[line width=5pt,red!70,opacity=0.35,line cap=round,line join=round] (0.825,1.050) -- (1.650,0.000) -- (3.300,0.000);
  \node[var,drp] (x0) at (0.000,0.000) {$\x_{t-1}$};
  \node[var,giv] (x1) at (1.650,0.000) {$\x_{t}$};
  \node[var,qry] (x2) at (3.300,0.000) {$\x_{t+1}$};
  \node[var] (x3) at (4.950,0.000) {$\x_{t+2}$};
  \node[var,drp] (z0) at (0.000,-1.250) {$\z_{t-1}$};
  \node[var,drp] (z1) at (1.650,-1.250) {$\z_{t}$};
  \node[var] (z2) at (3.300,-1.250) {$\z_{t+1}$};
  \node[var] (z3) at (4.950,-1.250) {$\z_{t+2}$};
  \node[ctl,drp] (u0) at (0.825,1.050) {$\uu_{t-1}$};
  \node[ctl,giv] (u1) at (2.475,1.050) {$\uu_{t}$};
  \node[ctl] (u2) at (4.125,1.050) {$\uu_{t+1}$};
  \draw[->,thick] (x0) -- (x1);
  \draw[->,thick] (x1) -- (x2);
  \draw[->,thick] (x2) -- (x3);
  \draw[->,thick] (u0) -- (x1);
  \draw[->,thick] (u1) -- (x2);
  \draw[->,thick] (u2) -- (x3);
  \draw[->,thick] (x0) -- (z0);
  \draw[->,thick] (x1) -- (z1);
  \draw[->,thick] (x2) -- (z2);
  \draw[->,thick] (x3) -- (z3);
  \draw[->,thick,dotted] ($(x0)+(-1.1,0)$) -- (x0);
  \draw[->,thick,dotted] (x3) -- ++(1.1,0);
\end{tikzpicture}

\end{center}

\emph{Check.} Each of the variables to drop is joined to the rest of the
graph only through $\x_t$. So every path from one of them to $\x_{t+1}$
passes through $\x_t$ and leaves it by the arrow
$\x_t\rightarrow\x_{t+1}$: there $\x_t$ is a chain node
($\x_{t-1}\rightarrow\x_t\rightarrow\x_{t+1}$) or a fork node
($\z_t\leftarrow\x_t\rightarrow\x_{t+1}$). It is given, so every path is
blocked.

\subsection*{(c) Prediction: the next control says nothing about the
current pose}

\emph{Claim.} $p(\x_t\mid\z_{1:t},\uu_{0:t})=p(\x_t\mid\z_{1:t},
\uu_{0:t-1})$: given the measurements and the earlier controls, $\x_t$ is
independent of $\uu_t$. This is the step of the proof of \cref{thm:bayes}
that turns the second factor into $\bel(\x_t)$.

\begin{center}
\begin{tikzpicture}[scale=1.25,>=Stealth,
  var/.style={circle,draw,thick,minimum size=9.5mm,inner sep=0pt,fill=white,font=\small},
  ctl/.style={rectangle,draw,minimum size=6.5mm,inner sep=1.5pt,fill=white,font=\small},
  giv/.style={fill=defblue!22},
  qry/.style={draw=orange!85!black,line width=1.6pt}]
  \draw[line width=5pt,red!70,opacity=0.35,line cap=round,line join=round] (1.650,0.000) -- (3.300,0.000) -- (2.475,1.050);
  \node[var] (x0) at (0.000,0.000) {$\x_{t-1}$};
  \node[var,qry] (x1) at (1.650,0.000) {$\x_{t}$};
  \node[var] (x2) at (3.300,0.000) {$\x_{t+1}$};
  \node[var] (x3) at (4.950,0.000) {$\x_{t+2}$};
  \node[var,giv] (z0) at (0.000,-1.250) {$\z_{t-1}$};
  \node[var,giv] (z1) at (1.650,-1.250) {$\z_{t}$};
  \node[var] (z2) at (3.300,-1.250) {$\z_{t+1}$};
  \node[var] (z3) at (4.950,-1.250) {$\z_{t+2}$};
  \node[ctl,giv] (u0) at (0.825,1.050) {$\uu_{t-1}$};
  \node[ctl,qry] (u1) at (2.475,1.050) {$\uu_{t}$};
  \node[ctl] (u2) at (4.125,1.050) {$\uu_{t+1}$};
  \draw[->,thick] (x0) -- (x1);
  \draw[->,thick] (x1) -- (x2);
  \draw[->,thick] (x2) -- (x3);
  \draw[->,thick] (u0) -- (x1);
  \draw[->,thick] (u1) -- (x2);
  \draw[->,thick] (u2) -- (x3);
  \draw[->,thick] (x0) -- (z0);
  \draw[->,thick] (x1) -- (z1);
  \draw[->,thick] (x2) -- (z2);
  \draw[->,thick] (x3) -- (z3);
  \draw[->,thick,dotted] ($(x0)+(-1.1,0)$) -- (x0);
  \draw[->,thick,dotted] (x3) -- ++(1.1,0);
\end{tikzpicture}

\end{center}

\emph{Check.} The only arrow at $\uu_t$ is $\uu_t\rightarrow\x_{t+1}$, so
every path from $\uu_t$ starts with it. To reach $\x_t$, the path must
then come back along $\x_{t+1}\leftarrow\x_t$: the two arrows point into
$\x_{t+1}$, which is a collider. Neither $\x_{t+1}$ nor any of its
descendants ($\z_{t+1}$, $\x_{t+2}$, \dots) is given, since only
measurements up to time $t$ are. The path is blocked. A path that
continues from $\x_{t+1}$ into the future cannot return to $\x_t$ without
passing through $\x_{t+1}$ a second time, so there is no other path.

The same argument shows that all the future controls are independent of
the past poses, which is used in the proof of \cref{prop:joint}.

\subsection*{(d) Smoothing: a later measurement does say something about
an earlier pose}

\emph{Question.} Given the measurements up to time $t$, is $\x_t$
independent of the next measurement $\z_{t+1}$?

\begin{center}
\begin{tikzpicture}[scale=1.25,>=Stealth,
  var/.style={circle,draw,thick,minimum size=9.5mm,inner sep=0pt,fill=white,font=\small},
  ctl/.style={rectangle,draw,minimum size=6.5mm,inner sep=1.5pt,fill=white,font=\small},
  giv/.style={fill=defblue!22},
  qry/.style={draw=orange!85!black,line width=1.6pt},
  drp/.style={draw=violet!80!black,line width=1.6pt}]
  \draw[line width=5pt,green!60!black,opacity=0.35,line cap=round,line join=round] (1.650,0.000) -- (3.300,0.000) -- (3.300,-1.250);
  \node[var] (x0) at (0.000,0.000) {$\x_{t-1}$};
  \node[var,qry] (x1) at (1.650,0.000) {$\x_{t}$};
  \node[var] (x2) at (3.300,0.000) {$\x_{t+1}$};
  \node[var] (x3) at (4.950,0.000) {$\x_{t+2}$};
  \node[var,giv] (z0) at (0.000,-1.250) {$\z_{t-1}$};
  \node[var,giv] (z1) at (1.650,-1.250) {$\z_{t}$};
  \node[var,drp] (z2) at (3.300,-1.250) {$\z_{t+1}$};
  \node[var] (z3) at (4.950,-1.250) {$\z_{t+2}$};
  \node[ctl,giv] (u0) at (0.825,1.050) {$\uu_{t-1}$};
  \node[ctl,giv] (u1) at (2.475,1.050) {$\uu_{t}$};
  \node[ctl,giv] (u2) at (4.125,1.050) {$\uu_{t+1}$};
  \draw[->,thick] (x0) -- (x1);
  \draw[->,thick] (x1) -- (x2);
  \draw[->,thick] (x2) -- (x3);
  \draw[->,thick] (u0) -- (x1);
  \draw[->,thick] (u1) -- (x2);
  \draw[->,thick] (u2) -- (x3);
  \draw[->,thick] (x0) -- (z0);
  \draw[->,thick] (x1) -- (z1);
  \draw[->,thick] (x2) -- (z2);
  \draw[->,thick] (x3) -- (z3);
  \draw[->,thick,dotted] ($(x0)+(-1.1,0)$) -- (x0);
  \draw[->,thick,dotted] (x3) -- ++(1.1,0);
\end{tikzpicture}

\end{center}

\emph{Check.} The path $\x_t\rightarrow\x_{t+1}\rightarrow\z_{t+1}$ is a
chain whose middle node $\x_{t+1}$ is not given. Nothing blocks it: the
path is open, and the graph does not give independence. Indeed $\z_{t+1}$
is a measurement of $\x_{t+1}$, which is one step away from $\x_t$. A
filter ignores this information about the past; a smoother uses it
(\cref{sec:smoothing}).

\subsection*{Complement: controls computed during the run}

In the network of \cref{fig:bayesnet} no arrow enters a control: the
controls are fixed in advance. A mower can also compute them as it goes,
for instance to steer back towards its planned path. What this changes
depends on what the controls are computed from. In the two drawings below
the new arrows are blue.

\emph{From the past measurements.} If $\uu_t$ is computed from
$\z_{1:t}$, arrows go from the measurements to the controls (only the
arrow from the latest measurement is drawn).

\begin{center}
\begin{tikzpicture}[scale=1.25,>=Stealth,
  var/.style={circle,draw,thick,minimum size=9.5mm,inner sep=0pt,fill=white,font=\small},
  ctl/.style={rectangle,draw,minimum size=6.5mm,inner sep=1.5pt,fill=white,font=\small},
  giv/.style={fill=defblue!22},
  qry/.style={draw=orange!85!black,line width=1.6pt}]
  \draw[line width=5pt,red!70,opacity=0.35,line cap=round,line join=round] (1.650,0.000) -- (3.300,0.000) -- (2.475,1.050);
  \draw[line width=5pt,red!70,opacity=0.35,line cap=round,line join=round] (1.650,0.000) -- (1.650,-1.250) to[out=35,in=-100] (2.475,1.050);
  \node[var] (x0) at (0.000,0.000) {$\x_{t-1}$};
  \node[var,qry] (x1) at (1.650,0.000) {$\x_{t}$};
  \node[var] (x2) at (3.300,0.000) {$\x_{t+1}$};
  \node[var] (x3) at (4.950,0.000) {$\x_{t+2}$};
  \node[var,giv] (z0) at (0.000,-1.250) {$\z_{t-1}$};
  \node[var,giv] (z1) at (1.650,-1.250) {$\z_{t}$};
  \node[var] (z2) at (3.300,-1.250) {$\z_{t+1}$};
  \node[var] (z3) at (4.950,-1.250) {$\z_{t+2}$};
  \node[ctl,giv] (u0) at (0.825,1.050) {$\uu_{t-1}$};
  \node[ctl,qry] (u1) at (2.475,1.050) {$\uu_{t}$};
  \node[ctl] (u2) at (4.125,1.050) {$\uu_{t+1}$};
  \draw[->,thick] (x0) -- (x1);
  \draw[->,thick] (x1) -- (x2);
  \draw[->,thick] (x2) -- (x3);
  \draw[->,thick] (u0) -- (x1);
  \draw[->,thick] (u1) -- (x2);
  \draw[->,thick] (u2) -- (x3);
  \draw[->,thick] (x0) -- (z0);
  \draw[->,thick] (x1) -- (z1);
  \draw[->,thick] (x2) -- (z2);
  \draw[->,thick] (x3) -- (z3);
  \draw[->,thick,dotted] ($(x0)+(-1.1,0)$) -- (x0);
  \draw[->,thick,dotted] (x3) -- ++(1.1,0);
  \draw[->,thick,blue!70!black] (z0) to[out=35,in=-100] (u0);
  \draw[->,thick,blue!70!black] (z1) to[out=35,in=-100] (u1);
  \draw[->,thick,blue!70!black] (z2) to[out=35,in=-100] (u2);
\end{tikzpicture}

\end{center}

Statement (c) still holds. There is a new path from $\uu_t$ to $\x_t$,
namely $\uu_t\leftarrow\z_t\leftarrow\x_t$, in which $\z_t$ is a chain
node. It is given, so the path is blocked; the same holds for a path
through an earlier measurement. The filter is therefore unchanged: it does
not matter how the controls were chosen, as long as they are computed from
measurements that the filter is given.

One thing does change. A \emph{later} control now carries information
about the past: the path $\uu_{t+1}\leftarrow\z_{t+1}\leftarrow\x_{t+1}$
is open as long as $\z_{t+1}$ is not given. The step of the proof of
\cref{prop:joint} that removes the later controls is then not valid. The
product \eqref{eq:joint} is still the joint density of the poses and the
measurements, with each $\uu_t$ replaced by its expression in terms of
$\z_{1:t}$, but it can no longer be read as a density given the whole
sequence of controls.

\emph{From the true pose.} If $\uu_t$ is computed from $\x_t$ itself,
there is an arrow $\x_t\rightarrow\uu_t$.

\begin{center}
\begin{tikzpicture}[scale=1.25,>=Stealth,
  var/.style={circle,draw,thick,minimum size=9.5mm,inner sep=0pt,fill=white,font=\small},
  ctl/.style={rectangle,draw,minimum size=6.5mm,inner sep=1.5pt,fill=white,font=\small},
  giv/.style={fill=defblue!22},
  qry/.style={draw=orange!85!black,line width=1.6pt},
  drp/.style={draw=violet!80!black,line width=1.6pt}]
  \draw[line width=5pt,red!70,opacity=0.35,line cap=round,line join=round] (1.650,0.000) -- (3.300,0.000) -- (2.475,1.050);
  \draw[line width=5pt,green!60!black,opacity=0.35,line cap=round,line join=round] (1.650,0.000) -- (2.475,1.050);
  \node[var] (x0) at (0.000,0.000) {$\x_{t-1}$};
  \node[var,qry] (x1) at (1.650,0.000) {$\x_{t}$};
  \node[var] (x2) at (3.300,0.000) {$\x_{t+1}$};
  \node[var] (x3) at (4.950,0.000) {$\x_{t+2}$};
  \node[var,giv] (z0) at (0.000,-1.250) {$\z_{t-1}$};
  \node[var,giv] (z1) at (1.650,-1.250) {$\z_{t}$};
  \node[var] (z2) at (3.300,-1.250) {$\z_{t+1}$};
  \node[var] (z3) at (4.950,-1.250) {$\z_{t+2}$};
  \node[ctl,giv] (u0) at (0.825,1.050) {$\uu_{t-1}$};
  \node[ctl,drp] (u1) at (2.475,1.050) {$\uu_{t}$};
  \node[ctl] (u2) at (4.125,1.050) {$\uu_{t+1}$};
  \draw[->,thick] (x0) -- (x1);
  \draw[->,thick] (x1) -- (x2);
  \draw[->,thick] (x2) -- (x3);
  \draw[->,thick] (u0) -- (x1);
  \draw[->,thick] (u1) -- (x2);
  \draw[->,thick] (u2) -- (x3);
  \draw[->,thick] (x0) -- (z0);
  \draw[->,thick] (x1) -- (z1);
  \draw[->,thick] (x2) -- (z2);
  \draw[->,thick] (x3) -- (z3);
  \draw[->,thick,dotted] ($(x0)+(-1.1,0)$) -- (x0);
  \draw[->,thick,dotted] (x3) -- ++(1.1,0);
  \draw[->,thick,blue!70!black] (x0) -- (u0);
  \draw[->,thick,blue!70!black] (x1) -- (u1);
  \draw[->,thick,blue!70!black] (x2) -- (u2);
\end{tikzpicture}

\end{center}

This arrow is a path between $\x_t$ and $\uu_t$ with no node in between,
so nothing can block it: statement (c) fails. The control is then a kind
of measurement of the pose, and a filter that drops $\uu_t$ in its
prediction step, as in the proof of \cref{thm:bayes}, leaves that
information unused. Assumptions \ref{A1} and \ref{A2} themselves still
hold, because the noises $\bepsilon_t$ and $\vv_t$ remain independent of
everything that came before.

\begin{mower}
The simulations of this book are in the second case at some steps. The
mower follows a path fixed in advance, but near the edge of the lawn the
pilot of the simulation, which knows the true pose, adds a turn towards
the centre. At those steps the commanded control says ``the mower is near
the edge'', and the filters do not use this: there they are not exactly
the Bayes filter of \cref{thm:bayes}, but a filter that ignores part of
what it is told.
\end{mower}
